\section{Experimental Setup and Metrics}\label{sec:setup}
\textbf{Hardware and software.} All training, evaluation and benchmarking ran locally on one NVIDIA GeForce
GTX~1660~SUPER (6\,GB GDDR6, Turing \texttt{sm\_75}, no tensor cores), an AMD Ryzen~5~2600 and 16\,GB RAM, under
Windows~11 with driver 621.14, Python~3.12, PyTorch~2.14.1 (CUDA~12.6), Transformers~4.57.6 and PEFT~0.21.2
(exact versions pinned in \texttt{uv.lock}). No paid service or cloud GPU was used. The application also runs on
CPU (slowly) and on the author's second machine (RTX~4050, 6\,GB).

\textbf{Decoding settings (all comparisons).} Greedy argmax decoding (no sampling, no repetition penalty), batch
size 1, at most 64 new tokens, bf16 weights, the project's own KV-cached decoder, and stopping at
\texttt{<|im\_end|>}. Base and fine-tuned models use the identical prompt, parser and examples. Benchmarks
disable EOS so that every run produces exactly the requested number of tokens.

\textbf{Protocol.} Model and epoch selection used only the validation split (8 patients, 206 annotations).
The test split (8 patients, 206 annotations) was decoded exactly once by \texttt{evaluation/final.py}, which first
writes a \texttt{FROZEN.json} record (adapter SHA-256, commit, prompt hash, decoding settings) and refuses to
run again; ordinary evaluation commands raise an error when asked for the test split.

\subsection{Metrics}\label{sec:metrics}
\textbf{Criterion classification.} Accuracy; macro-F1 over all six labels (a class absent from the gold
labels and never predicted scores $F_1=0$); macro-F1$_p$ over the labels present in the gold set of that
split (needed because the test split has no \emph{excluded} example); per-class precision, recall and F1;
and confusion matrices. Everything is reported overall and separately for inclusion criteria (4 labels)
and exclusion criteria (4 labels). Unparseable outputs or labels invalid for the criterion type are a
separate \emph{invalid} prediction that always counts as an error. Because the 8 test patients are the
independent units, we report 95\% patient-level bootstrap intervals (1{,}000 resamples).

\textbf{Evidence localisation.} For each record with gold set $G$ and predicted set $P$ of sentence IDs:
$P{=}G{=}\emptyset$ gives precision = recall = F1 = 1 (correct abstention); $G{=}\emptyset\neq P$ gives
precision 0, recall 1, F1 0 (hallucinated evidence); $G\neq\emptyset{=}P$ gives precision 1, recall 0, F1 0
(missed evidence); otherwise standard set precision/recall/F1. \emph{Micro} scores pool TP/FP/FN over all
records (empty--empty records contribute nothing); \emph{macro} scores average the per-record values. Because
conventions could inflate scores, we also report micro-F1 on the subset with non-empty gold evidence, and
count correct abstentions, hallucinations and cited IDs that do not exist in the note.

\textbf{Inference benchmarks.} Same model and inputs for every compared path; warm-up runs first (one for
decoding, two for attention), then repeated timed runs (three for KV-cache, two per prompt for speculative
decoding, five for attention) with \texttt{torch.cuda.synchronize()} around each timed region;
peak memory is \texttt{torch.cuda.max\_memory\_allocated()} after \texttt{reset\_peak\_memory\_stats()},
reported as extra memory over the resident weights. KV-cache inputs of 128/512/1024 tokens are built from real
validation prompts (several notes concatenated for the longest), with 16/64/128 output tokens (the quadratic no-cache reference makes longer outputs impractical: 128 tokens after a 1{,}024-token input already take about four minutes per run).
